%! Parent Visit — Unit 2 — "Five Ways to Factor" — one-page handout
% STANDALONE handout for the parent-visit class. Merged into no packet — see ../Makefile.
% ONE PAGE. No name row: it goes home with the parent. No separate key: the answers are below
% for the teacher, and each parent checks their student's board work by multiplying back out.
% Five students, five methods from Lesson 2.2 (one problem each, at the board, side by side):
%   1  GCF                     8x^3-12x^2+20x  = 4x(2x^2-3x+5)        (2x^2-3x+5 is prime: D<0)
%   2  Difference of squares   4x^2-81         = (2x+9)(2x-9)
%   3  Trinomial, a = 1        x^2-2x-35       = (x-7)(x+5)
%   4  Grouping (master prod.) 6x^2+11x-10     : ac=-60 -> 15,-4; 6x^2+15x-4x-10 = 3x(2x+5)-2(2x+5)
%                                              = (3x-2)(2x+5)
%   5  Slip and slide          4x^2-4x-15      : slip x^2-4x-60 = (x-10)(x+6); slide
%                                              (x-10/4)(x+6/4) = (x-5/2)(x+3/2); up (2x-5)(2x+3)
% The run (about 10 min):
%   0-1  Five students go to the board, one problem each; parents read the method box.
%   1-5  Students work silently. Parents write each answer in the blank as it appears.
%   5-9  Each student narrates one step in one sentence ("I slipped the 4 onto the 15 ...").
%        Parents check one answer by multiplying it back out (the Check box).
%   9-10 The crux: ask a parent why #1 starts with the GCF; a student answers that #4 and #5
%        both fail if a GCF is left in.
\documentclass[12pt]{article}
\usepackage{algebra2-article}
\usepackage{algebra2-boxes}

\begin{document}

\pageheader{Unit 2: Quadratic Functions}{Parent Visit --- Five Ways to Factor}

\boxguard
\begin{scenariobox}[Tonight at the Board]{goldacc}
\textbf{Factoring} rewrites a sum as a \textbf{product} --- multiplying, run backwards. It is how
we \emph{solve} quadratics: if $(2x-5)(2x+3)=0$, one factor must be zero. Five students each
factor one expression at the board, \textbf{each by a different method}. Copy their answers below.
\end{scenariobox}

\vspace{0.06in}
\boxguard[24]
\begin{scenariobox}[The Five Problems]{forest}
\small
\renewcommand{\arraystretch}{1.55}
\begin{tabularx}{\linewidth}{@{}c >{\raggedright\arraybackslash}p{2.6cm} >{\raggedright\arraybackslash}X >{\raggedright\arraybackslash}p{2.7cm} l@{}}
& \textbf{Method} & \textbf{The move} & \textbf{Problem} & \textbf{Answer} \\ \hline
\textbf{1} & GCF
  & Pull out the largest factor every term shares. \textbf{Always first.}
  & $8x^3-12x^2+20x$ & \blank{2.9cm} \\
\textbf{2} & Diff.\ of squares
  & $A^2-B^2=(A+B)(A-B)$.
  & $4x^2-81$ & \blank{2.9cm} \\
\textbf{3} & Trinomial $a=1$
  & Two numbers that \emph{multiply} to $c$ and \emph{add} to $b$.
  & $x^2-2x-35$ & \blank{2.9cm} \\
\textbf{4} & Grouping
  & Master product $a\cdot c$: its pair splits $bx$; GCF each half.
  & $6x^2+11x-10$ & \blank{2.9cm} \\
\textbf{5} & Slip and slide
  & Slip $a$ onto $c$, factor, slide $a$ back, reduce, slide up.
  & $4x^2-4x-15$ & \blank{2.9cm} \\
\end{tabularx}
\end{scenariobox}

\vspace{0.06in}
\boxguard[12]
\begin{scenariobox}[Same Problem, Two Roads: \#4 and \#5 Both Start With $a\cdot c$]{navy}
\small
\begin{minipage}[t]{0.47\linewidth}\vspace{0pt}\raggedright
\textbf{Grouping:} $3x^2+10x+8$\par\vspace{1pt}
$a\cdot c=24$; the pair is $6$ and $4$.\par
$3x^2+6x+4x+8$\par
$3x(x+2)+4(x+2)$\par
$(3x+4)(x+2)$ --- the brackets match.
\end{minipage}\hfill
\begin{minipage}[t]{0.49\linewidth}\vspace{0pt}\raggedright
\textbf{Slip and slide:} $3x^2+10x+8$\par\vspace{1pt}
Slip: $x^2+10x+24=(x+6)(x+4)$\par
Slide: $\left(x+\frac{6}{3}\right)\left(x+\frac{4}{3}\right)$\par
Reduce: $(x+2)\left(x+\frac{4}{3}\right)$\par
Slide up: $(x+2)(3x+4)$ --- same answer.
\end{minipage}
\end{scenariobox}

\vspace{0.06in}
\boxguard[5]
\begin{scenariobox}[Parent: Check One Answer]{slate}
Pick one board answer and \textbf{multiply it back out}. Do you get the original?\par
\vspace{6pt}
Problem \# \blank{1.0cm} \qquad Multiplied out: \blank{6.2cm} \qquad $\checkmark$ \; or \; $\times$
\end{scenariobox}

\vspace{0.06in}
\begin{remindbox}
\textbf{For parents.} Unit~2 is quadratics: graphing parabolas, then solving them by factoring,
square roots, and the quadratic formula. Homework is due the first class after two study halls. \textbf{Ask your student:}
\emph{``Why do you always look for a GCF first?''}
\end{remindbox}

\end{document}
